\documentclass[11pt]{article}
\usepackage[margin=1in]{geometry}

% Core packages
\usepackage{amsmath,amssymb}
\usepackage{tikz-cd}
\usepackage{multicol}
\usepackage{booktabs}
\usepackage{longtable}
\usepackage{array}
\usepackage{xcolor}

\newcommand{\modelbetter}{\textcolor{green!50!black}{\checkmark}}
\newcommand{\modelnotbetter}{\textcolor{red!75!black}{\ensuremath{\times}}}
\newcommand{\modelconditional}{\textcolor{orange!80!black}{\ensuremath{\triangle}}}
\newcommand{\modeluntested}{\textcolor{gray}{---}}

% Paragraphs
\setlength{\parindent}{0pt}
\setlength{\parskip}{1\baselineskip}

\title{Patient Representations from Bulk RNA-seq:
From Gene Expression to Foundation Models}
\date{\today}

\begin{document}
\maketitle
\begingroup
\setlength{\parskip}{0pt}
\tableofcontents
\endgroup

\section{Introduction}

Biological patient representations turn measurements from one or more
modalities into a form that a predictive model can use. This is often a single
vector, although some methods produce a set or sequence of vectors that is
later aggregated. A useful representation should retain the biological signal
needed for tasks such as predicting prognosis or treatment response.

Constructing such representations is difficult, and many approaches are
reasonable. This post is a short, curated literature review of representations
derived from bulk gene expression, with particular attention to pathways and
gene modules. These structures provide useful biological inductive biases and
raise interesting modeling questions.

Although the ideas behind deep learning for bulk RNA are interesting, the
available evidence does not show that it consistently outperforms simpler
representations such as normalized gene expression or PCA. The problem remains
important, however, and progress depends partly on understanding and combining
ideas from previous work.

This is the post I wish had been available when I began reading about bulk RNA
modeling, and I hope it will be useful to others. A future follow-up may focus
on histology.

I have likely missed some valuable studies and ideas; contributions and pull
requests are welcome.

\section{Bulk gene-expression-based representations}

\subsection{Gene expression as the representation}

For $N$ patients and $G$ measured genes, write the gene-expression data as
\begin{equation}
    X = (x_{ig}) \in \mathbb{R}^{N \times G}.
\end{equation}
Here, $x_{ig}$ is the expression of gene $g$ in patient $i$. The simplest
patient representation is therefore the corresponding row $x_i\in
\mathbb{R}^{G}$ after appropriate preprocessing. For RNA-seq, common choices
include CPM or TPM normalization and transformations such as
$\log(\mathrm{TPM}+1)$.

For example, Golub et al. used DNA-microarray gene-expression profiles as
patient representations to distinguish acute myeloid leukemia (AML) from acute
lymphoblastic leukemia (ALL) and to construct a predictor for new leukemia
samples \cite{golub1999molecular}. Their classifier used preprocessed
expression values after selecting an informative subset of genes. The important
point is that gene expression itself can already be highly predictive.

This continues well into the current era. A 2024 benchmark of bulk RNA-seq
representations found that using the transcript-level expression vector itself
(``Identity'') or PCA matched or outperformed more complex deep-learning
representations for cancer-survival prediction \cite{gross2024robust}. We expand on this at the end of this post.

\subsection{Finding data-driven gene modules with WGCNA}

Weighted Gene Co-expression Network Analysis (WGCNA) groups genes that vary
together across patients \cite{langfelder2008wgcna}. It first constructs a
soft-thresholded adjacency matrix by raising each pairwise correlation to a
power $\beta$:
\begin{equation}
    \begin{aligned}
        \rho_{gh} &= \operatorname{cor}(x_{\cdot g},x_{\cdot h}), \\
        a_{gh} &= |\rho_{gh}|^{\beta}.
    \end{aligned}
\end{equation}
Strongly correlated genes receive high edge weights, while weak correlations
are suppressed.

WGCNA uses topological overlap to compare shared neighborhoods, then applies
hierarchical clustering to identify gene modules.

A key distinction is that WGCNA constructs modules from the cohort itself.
An alternative is to use predefined gene sets. MSigDB is a broad collection of
annotated gene sets assembled from pathways, ontologies, and published
expression signatures \cite{liberzon2011msigdb}. Reactome provides manually
curated biological reactions and pathways \cite{milacic2024reactome}, while
KEGG provides pathway maps linking genes, molecules, and cellular processes
\cite{kanehisa2000kegg}. These standardized gene sets are easier to compare
across studies, although they may miss structure specific to a particular
cohort.


\subsection{Summarizing gene modules and pathways}

\subsubsection{PC1 scores}

WGCNA summarizes each module by the first principal component (PC1) of its gene
expression matrix. This summary is called the module eigengene. Concatenating
the module eigengene scores gives one vector per patient.

One subtlety is that the sign of a PCA score is arbitrary. For a
module-expression
matrix $X_m$, PCA finds the direction
\begin{equation}
    w_1 = \underset{\lVert w\rVert_2=1}{\operatorname{arg\,max}}
    \;\operatorname{Var}(X_m w).
\end{equation}
Both $w_1$ and $-w_1$ solve this problem, so the sign of PC1 is arbitrary.
WGCNA orients the eigengene to correlate positively with the module's average
expression profile. After this orientation, positive and negative scores mean
that a patient lies above or below the cohort average along the module's shared
expression pattern.

\subsubsection{Single-sample Gene Set Enrichment Analysis (ssGSEA)}

Before describing ssGSEA, it is useful to note two properties of PC1 scores.

\begin{itemize}
    \item PC1 captures variation across patients. Genes may be highly expressed
    in every patient, but this will not produce high PC1 scores if their
    expression changes little across the cohort.
    \item PC1 does not ask whether a pathway's genes rank highly relative to
    other genes in the same patient. It places each patient along the dominant
    pattern of variation within the selected genes.
\end{itemize}

ssGSEA takes a different perspective: rather than asking how a pathway varies
across patients, it asks where the pathway's genes appear in each individual
patient's ranked transcriptome \cite{barbie2009systematic}.

For one patient, ssGSEA ranks all genes by expression. Walking down this list,
it increases a running score when it encounters a pathway gene and decreases
the score otherwise.

Let $D(k)$ denote the running enrichment score after encountering the gene at
position $k$ in the ranked list. It is updated as
\begin{equation}
    D(k) = D(k-1) + \Delta_k,
\end{equation}
where, intuitively,
\begin{equation}
    \Delta_k =
    \begin{cases}
        +\text{rank-weighted contribution}, & \text{if the gene belongs to the pathway}, \\[4pt]
        -\text{constant contribution}, & \text{otherwise}.
    \end{cases}
\end{equation}
More precisely, for a pathway gene set $S$,
\begin{equation}
    \Delta_k =
    \begin{cases}
        \displaystyle
        \frac{r_k^{\alpha}}
        {\sum_{j:g_j \in S} r_j^{\alpha}},
        & g_k \in S, \\[10pt]
        \displaystyle
        -\frac{1}{G-|S|},
        & g_k \notin S,
    \end{cases}
\end{equation}
where $G$ is the total number of genes, $|S|$ is the number of genes in the
pathway, $r_k$ is a positive weight derived from the rank of gene $g_k$, and
$\alpha$ controls how strongly the highest-ranked genes are weighted. The
ssGSEA score sums the running difference across the full ranking:
\begin{equation}
    \operatorname{ES}_{\mathrm{ssGSEA}}(S) = \sum_{k=1}^{G} D(k).
\end{equation}
Thus PC1 is a cohort-level summary of how a set of genes varies together,
whereas ssGSEA is computed within one patient and measures where a pathway's
genes fall in that patient's ranked transcriptome.

These summaries are worth considering for two reasons. First, they are strong,
interpretable baselines for bulk RNA-seq and are easy to overlook when comparing
modern architectures. Second, the same ideas can be generalized into learned
models, as the following sections illustrate.


\subsection{Learned pathway embeddings}

Instead of summarizing a pathway with a fixed statistic, we can learn a
pathway-specific representation. SurvPath provides one example
\cite{jaume2024modeling}. For pathway $\mathcal{P}_m$, let
$x_{\mathcal{P}_m}\in\mathbb{R}^{|\mathcal{P}_m|}$ contain the expression
values of its member genes. The pathway embedding is
\begin{equation}
    z_m = \phi_m\!\left(x_{\mathcal{P}_m}\right) \in \mathbb{R}^{d},
\end{equation}
where $\phi_m$ is a two-layer neural network specific to pathway
$\mathcal{P}_m$. Because pathways contain different numbers of genes, these
networks have different input dimensions and separate parameters, but all
produce embeddings of the same dimension $d$.

SurvPath stacks these pathway embeddings into a sequence and combines them
with histology-patch embeddings using a multimodal Transformer. The pathway
encoders and Transformer are trained end to end for survival prediction. This
encourages each pathway embedding to retain task-relevant information, which
the Transformer aggregates into a patient-level prediction.

\subsection{Learned gene--pathway graph embeddings}

The independent encoders above do not account for the fact that pathways
overlap. A gene may belong to several pathways, but each copy is processed in
isolation. ProtoPathway is a recent work that instead represents genes and pathways as the two node
types of a bipartite graph \cite{gallaghersyed2026protopathway}. A gene $g_i$
is connected to a pathway $p_j$ exactly when it belongs to that pathway:
\begin{equation}
    (g_i,p_j)\in E \quad\Longleftrightarrow\quad g_i\in p_j.
\end{equation}
For each patient, gene nodes are initialized with their expression values and
pathway nodes with zeros.

The first layers use GraphSAGE with mean aggregation
\cite{hamilton2017inductive}. In simplified form,
\begin{equation}
    m_v^{(l)} = \frac{1}{|\mathcal{N}(v)|}
    \sum_{u\in\mathcal{N}(v)}h_u^{(l)},
    \qquad
    h_v^{(l+1)} = \sigma\!\left(
    W^{(l)}[h_v^{(l)}\mathbin\Vert m_v^{(l)}]
    \right).
\end{equation}
Messages travel in both directions. A gene shared by two pathways therefore
provides a route through which those pathways can exchange information. The
resulting pathway vectors describe both their member genes and their context in
the wider pathway graph. Unlike the independent encoders above, the GNN shares
parameters across all pathways.

A final GATv2 layer learns how much each neighboring gene should contribute to
a pathway embedding:
\begin{equation}
    z_p = \sum_{g\in\mathcal{N}(p)} \alpha_{gp} W h_g.
\end{equation}
GATv2 makes these weights depend on the pathway doing the querying
\cite{brody2022attentive}. A gene shared by two pathways can consequently have
different weights in each, so $\alpha_{g,p_1}$ need not equal
$\alpha_{g,p_2}$. ProtoPathway therefore retains named pathway representations
while allowing overlapping pathways to contextualize one another through their
shared genes.


\subsection{Pathway crosstalk with BinoX}

One useful idea is to estimate pathway-to-pathway relationships from a
gene-level interaction network. BinoX provides such a prior for the Pathformer
attention mechanism described next.

BinoX estimates whether two pathways are more functionally connected than
expected by chance \cite{ogris2017binox}. It uses a functional association
network such as FunCoup \cite{schmitt2014funcoup}. Its evidence comes mainly
from public databases and experimental datasets. FunCoup combined evidence
from interactions, expression, regulation, cellular localization, protein
domains, and evolutionary relationships.
Many of these resources curate published experiments, so literature contributes
indirectly. FunCoup scores each evidence type and combines the scores with a
Bayesian model. Its human network contained 18,113 genes and about 4.48 million
scored associations. These edges indicate functional evidence, not necessarily
direct physical interactions.

For two pathways $P_A$ and $P_B$, BinoX counts the FunCoup edges connecting
their genes:
\begin{equation}
    k = \sum_{i\in P_A}\sum_{j\in P_B}x_{ij},
\end{equation}
where $x_{ij}=1$ if genes $i$ and $j$ are connected and $x_{ij}=0$ otherwise.
Large pathways and pathways containing hub genes naturally have more edges, so
BinoX compares $k$ with degree-preserving randomized networks. It models the
null count and tests its upper tail:
\begin{equation}
    K\sim\operatorname{Binomial}(n,p),
    \qquad
    p_{\mathrm{cross}}=\Pr(K\geq k_{\mathrm{observed}}),
\end{equation}
where $p$ is estimated from the randomized networks. A small
$p_{\mathrm{cross}}$ indicates more pathway crosstalk than expected from the
background network.

Applying BinoX to every pathway pair produces a pathway-by-pathway crosstalk
matrix. This is fixed external biological knowledge rather than a
patient-specific quantity. The next section explains how Pathformer uses it.

\subsection{Contextualizing pathway embeddings with Pathformer}

Pathformer \cite{liu2024pathformer} creates a patient representation from
learned pathway embeddings. The full model uses ``criss-cross'' attention over both
pathways and omics modalities; to stay focused on bulk gene expression, we omit
the modality axis here.

First, a single pathway-based sparse neural network uses pathway membership to
restrict the allowed gene-to-pathway connections. In the RNA-only
configuration, it produces one scalar activity for each pathway:
\begin{equation}
    s_{ip}=f_{\theta,p}\!\left(x_{i,P_p}\right)\in\mathbb{R},
\end{equation}
where $x_{i,P_p}$ contains patient $i$'s expression values for genes assigned
to pathway $P_p$. The scalar is then multiplied by a learned
pathway-specific vector $e_p\in\mathbb{R}^{d}$ to form the Transformer token:
\begin{equation}
    z_{ip}=s_{ip}e_p\in\mathbb{R}^{d}.
\end{equation}
Thus the sparse network supplies the patient-specific pathway magnitude,
whereas $e_p$ supplies a learned pathway identity and direction in embedding
space. The multiplication is a particular, restrictive design choice rather
than a necessity. A natural alternative would be a pathway-specific MLP that
directly outputs a vector,
\begin{equation}
    z_{ip}=\phi_p\!\left(x_{i,P_p}\right)\in\mathbb{R}^{d},
\end{equation}
as in SurvPath. Whether Pathformer's scalar-times-vector construction provides
useful regularization or discards useful within-pathway information is an
interesting modeling question that could be verified by ablations. In the full multi-omics configuration, the
embedding step can instead be bypassed and the pathway-by-modality values are
passed directly to criss-cross attention. Pathformer uses 1,497 pathways
curated from KEGG, Reactome, PID, and BioCarta.

The BinoX analysis supplies Pathformer's initial pathway adjacency matrix,
denoted in the paper by $P^{(0)}\in\mathbb{R}^{N_p\times N_p}$, where
$N_p=1497$. To distinguish this matrix from the pathway sets, we write
\begin{equation}
    B_{\mathrm{BinoX}} = P^{(0)},
    \qquad
    (B_{\mathrm{BinoX}})_{ij}
    = s_{\mathrm{BinoX}}(P_i,P_j),
\end{equation}
where $s_{\mathrm{BinoX}}(P_i,P_j)$ is the BinoX-derived adjacency value for
pathways $P_i$ and $P_j$. This initial matrix is fixed before training and is
the same for every patient \cite{liu2024pathformer}.

The pathway embeddings are stacked into a matrix $Z$. In the first Transformer
block, attention combines patient-specific pathway similarity with the fixed
BinoX prior:
\begin{equation}
    Q=ZW_Q,\qquad K=ZW_K,\qquad V=ZW_V,
\end{equation}
\begin{equation}
    A^{(0)}=\operatorname{softmax}\!\left(
        \frac{QK^{\top}}{\sqrt{d}}+B_{\mathrm{BinoX}}
    \right),
    \qquad
    \widetilde{Z}^{(1)}=A^{(0)}V.
\end{equation}
Here, $QK^{\top}$ is computed from the current patient's pathway embeddings,
whereas $B_{\mathrm{BinoX}}$ supplies external biological knowledge.
Because it is added to the attention logits, BinoX acts as a bias rather than a
hard mask.

In later blocks, the crosstalk matrix is updated from similarities between that
patient's contextualized pathway embeddings:
\begin{equation}
    B_{\mathrm{BinoX}}
    \longrightarrow \text{Block 1}
    \longrightarrow B^{(1)}
    \longrightarrow \text{Block 2}
    \longrightarrow B^{(2)}
    \longrightarrow \text{Block 3}.
\end{equation}
Thus $B^{(1)}$ and later matrices are sample-specific; they are not correlations
computed across patients or batches. At these later stages, both $QK^{\top}$
and $B^{(l)}$ describe relationships between learned pathway embeddings. The
most distinctive role of the crosstalk term is therefore in the first block,
where BinoX introduces information that did not come from the patient data.

Pathformer uses biological knowledge twice: pathway membership constrains the
gene-to-pathway encoder, and BinoX initializes pathway-to-pathway attention.
It was evaluated on TCGA survival-risk, cancer-stage, and drug-response
classification tasks using multi-omics data.

\subsection{Gene-level foundation modeling with BulkFormer}

The methods above explicitly construct pathway representations. BulkFormer
instead operates directly on genes and uses a gene--gene graph as biological
prior knowledge \cite{kang2026bulkformer}. It covers 20,010 protein-coding
genes and was pretrained on 581,503 human bulk RNA-seq profiles.

BulkFormer constructs one input token for each gene. Its three
$d$-dimensional components are combined by element-wise addition, not
concatenation:
\begin{equation}
    h_{ig}^{(0)}
    = e_g^{\mathrm{ESM2}}
    + e_{\mathrm{REE}}(x_{ig}^{\mathrm{in}})
    + e_i^{\mathrm{sample}},
    \qquad
    e_i^{\mathrm{sample}}
    = \operatorname{MLP}(x_i^{\mathrm{in}}).
\end{equation}
Here, $e_g^{\mathrm{ESM2}}$ represents gene identity using a precomputed
embedding of the gene's encoded protein \cite{lin2023evolutionary}, rather than
a freely learned gene-ID lookup. The second term represents that gene's
continuous expression value. The final term is a compressed summary of the
sample-wide input expression vector, and the same vector is added to every gene
token. Thus each token contains the gene's identity, its own expression, and
global context from the sample. During masked pretraining,
$x_i^{\mathrm{in}}$ is the partially masked expression profile, so the sample
embedding does not directly include the held-out values.

BulkFormer encodes the continuous expression component with a Rotary Expression
Embedding (REE) \cite{kang2026bulkformer}. It maps an expression value $x_g$
across several frequencies:
\begin{equation}
    e_{\mathrm{REE}}(x_g)
    = \bigl(
        \sin(\omega_1x_g),\cos(\omega_1x_g),\ldots,
        \sin(\omega_{d/2}x_g),\cos(\omega_{d/2}x_g)
      \bigr).
\end{equation}
This deterministic, multi-frequency encoding preserves the continuity of gene
expression without first placing values into discrete bins. It adapts the
sinusoidal idea behind rotary position encodings, but uses expression magnitude
rather than token position.
This is an interesting design choice rather than a necessity: a small MLP could
also map each scalar expression value to a vector. REE instead supplies a
fixed, smooth basis at several scales, so nearby expression values have related
representations without requiring the model to learn that structure from
scratch. An MLP would be more flexible, but would introduce additional
parameters and could learn an irregular encoding, particularly in expression
ranges that are sparsely represented during training.

BulkFormer uses a hybrid GCN--Performer architecture. The GCN propagates
information along a predefined gene graph:
\begin{equation}
    H^{(l+1)} = \sigma\!\left(
        \widetilde{D}^{-1/2}\widetilde{A}
        \widetilde{D}^{-1/2}H^{(l)}W^{(l)}
    \right).
\end{equation}
Equivalently, the update for gene $g$ is
\begin{equation}
    h_g^{(l+1)} = \sigma\!\left(
        \sum_{u\in\mathcal{N}(g)\cup\{g\}}
        \frac{\widetilde{a}_{gu}}
        {\sqrt{\widetilde{d}_g\widetilde{d}_u}}
        h_u^{(l)}W^{(l)}
    \right),
\end{equation}
where $\widetilde{a}_{gu}$ is an entry of the adjacency matrix with self-loops
and $\widetilde{d}_g$ is its corresponding degree. Thus each gene
receives a degree-normalized weighted sum of its own representation and those
of its graph neighbors. The self-loop preserves the gene's own identity and
expression instead of replacing them with neighboring information. Degree
normalization is especially useful in biological graphs, where hub genes can
have many neighbors and would otherwise dominate the update.

The authors compared protein-interaction, Gene Ontology, and gene
co-expression graphs. The co-expression graph performed best and was used in
the final model. It connects each gene to at most its 20 strongest neighbors by
absolute Pearson correlation, discarding correlations below 0.4.

The Performer component lets genes exchange information globally. Ordinary
self-attention constructs a $G\times G$ matrix and therefore scales
quadratically with the number of genes:
\begin{equation}
    \operatorname{Attention}(Q,K,V)
    = \operatorname{softmax}\!\left(
        \frac{QK^{\top}}{\sqrt{d}}
    \right)V.
\end{equation}
Performer approximates the softmax kernel using a random feature map $\phi$,
\begin{equation}
    \exp\!\left(\frac{q^{\top}k}{\sqrt{d}}\right)
    \approx\phi(q)^{\top}\phi(k),
\end{equation}
which allows attention to be rearranged as
\begin{equation}
    \operatorname{Attention}(Q,K,V)
    \approx D^{-1}\phi(Q)\bigl(\phi(K)^{\top}V\bigr).
\end{equation}
Here, $D$ supplies the corresponding normalization. This avoids explicitly
constructing the $G\times G$ attention matrix and gives approximately linear
scaling in the number of genes \cite{choromanski2021rethinking}.

BulkFormer is pretrained by masking 15\% of expression values and reconstructing
their continuous values from the remaining transcriptome:
\begin{equation}
    \mathcal{L}_{\mathrm{mask}}
    = \frac{1}{|\mathcal{M}|}
    \sum_{g\in\mathcal{M}}(x_g-\widehat{x}_g)^2.
\end{equation}
Unlike the pathway-based methods above, BulkFormer learns contextualized gene
representations rather than explicit pathway representations. A natural
extension would add pathway tokens or a gene--pathway graph, combining
large-scale self-supervised pretraining with named pathway representations.

\section{Epilogue: Do learned representations improve on expression itself?}

The evidence is mixed and strongly task-dependent. Across several standardized
benchmarks, deep representations have not consistently improved on normalized
expression or PCA, particularly for survival prediction in modest-sized patient
cohorts. This does not imply that deep learning cannot help; rather, it shows
that additional model complexity is not itself sufficient. Table~\ref{tab:bulk-evaluations}
summarizes representative evaluations.

The final column records whether a study reported a direct advantage over an
expression or PCA baseline under its own evaluation protocol. The papers use
different endpoints and statistical criteria, so the symbols should not be
read as a common meta-analytic significance test. A green check indicates a
clear direct gain, a red cross indicates that expression or PCA was on par or
better, an orange triangle indicates a method- or setting-dependent gain, and
a dash indicates that a matched comparison was not reported. This distinction
matters for Pathformer and BulkFormer: both improved on their selected
baselines, but neither paper established superiority over a matched RNA-only
expression or PCA baseline across its tasks.

\begingroup
\footnotesize
\setlength{\parskip}{0pt}
\setlength{\tabcolsep}{4pt}
\renewcommand{\arraystretch}{1.2}
\begin{longtable}{
    >{\raggedright\arraybackslash}p{0.16\textwidth}
    >{\raggedright\arraybackslash}p{0.20\textwidth}
    >{\raggedright\arraybackslash}p{0.18\textwidth}
    >{\raggedright\arraybackslash}p{0.26\textwidth}
    >{\centering\arraybackslash}p{0.11\textwidth}}
\caption{Representative evaluations of bulk-transcriptomic modeling. Results
are not directly comparable across rows.}
\label{tab:bulk-evaluations}\\
\toprule
Study & Dataset & Task & Modeling method & Direct gain? \\
\midrule
\endfirsthead
\toprule
Study & Dataset & Task & Modeling method & Direct gain? \\
\midrule
\endhead
\bottomrule
\endfoot

Smith et al.\ \cite{smith2020standard}
& ${\sim}45{,}000$ recount2 samples
& 24 classification; 26 survival tasks
& No embedding, PCA, SDAE, and VAE; matched downstream predictors
& \modelnotbetter \\

Way and Greene \cite{way2017evaluating}
& ${>}10{,}000$ TCGA tumours
& NF1-inactivation prediction; HGSC subtype analysis
& Three VAEs versus PCA, ICA, NMF, and ADAGE
& \modelnotbetter \\

Gross et al.\ \cite{gross2024robust}
& TCGA: 11 cohorts and 33-cancer pan-cancer data
& Censoring-aware survival
& AE variants, VAE, MAE, and GNN versus Identity and PCA
& \modelnotbetter \\

Gross et al.\ \cite{gross2024robust}
& DepMap cell lines
& Gene-essentiality prediction
& The same representation families
& \modelconditional \\

Pathformer \cite{liu2024pathformer}
& TCGA multi-omics cohorts
& Risk group, stage, and drug response
& Pathway-informed Transformer versus 18 integration methods
& \modeluntested \\

GexBERT \cite{jiang2025gexbert}
& TCGA; 14 cancer-specific survival tasks
& Survival representation
& Transformer summary embedding versus 200-component PCA
& \modelbetter \\

GexBERT \cite{jiang2025gexbert}
& TCGA; 64--1,024 randomly selected input genes
& Survival from incomplete panels
& Observed plus restored anchor genes versus observed genes alone
& \modelconditional \\

Xia et al.\ \cite{xia2022crossstudy}
& NCI-60, CTRP, GDSC, CCLE, and gCSI
& Cross-study drug response
& RF, LightGBM, and deep neural networks, including UnoMT
& \modeluntested \\

SurvBoard \cite{wissel2025survboard}
& 28 cohorts from TCGA, ICGC, TARGET, and METABRIC
& Gene-expression-only survival
& Elastic net and survival RF versus neural-network models
& \modelnotbetter \\

BulkFormer \cite{kang2026bulkformer}
& ${>}500{,}000$ training profiles; multiple bulk RNA-seq benchmarks
& Five downstream tasks
& BulkFormer versus seven single-cell foundation models transferred to bulk data
& \modeluntested \\
\end{longtable}
\endgroup

\modelbetter{} reported direct gain; \modelnotbetter{} expression or PCA on par
or better; \modelconditional{} gain limited to some methods or input regimes;
\modeluntested{} no matched expression/PCA comparison.

For Gross et al., deep representations had a slight, consistent advantage on
the DepMap task overall, but some were equivalent rather than superior to the
baselines under the paper's pairwise acceptance criterion
\cite{gross2024robust}.

For GexBERT, the green check records the numerical improvement of its summary
embedding over PCA in the reported survival results, rather than a standardized
significance test. Reconstructed expression was useful mainly for small,
randomly selected panels; its benefit diminished for larger panels and was
negligible when the observed genes were already selected for prognostic
relevance \cite{jiang2025gexbert}.

The open research problem is therefore not simply to build a more expressive
model. A useful representation must retain patient-specific gene information,
exploit sufficiently large training cohorts, and introduce biological structure
without discarding predictive signal. Pathways and gene modules may provide
useful regularization and interpretability; hybrid representations that
preserve gene-level information remain a promising direction.



\begin{thebibliography}{99}

\bibitem{golub1999molecular}
T.~R. Golub, D.~K. Slonim, P. Tamayo, et al.
\newblock Molecular classification of cancer: Class discovery and class
prediction by gene expression monitoring.
\newblock \emph{Science}, 286(5439):531--537, 1999.
\newblock doi:10.1126/science.286.5439.531.

\bibitem{gross2024robust}
B. Gross, A. Dauvin, V. Cabeli, et al.
\newblock Robust evaluation of deep learning-based representation methods for
survival and gene essentiality prediction on bulk RNA-seq data.
\newblock \emph{Scientific Reports}, 14:17064, 2024.
\newblock doi:10.1038/s41598-024-67023-8.

\bibitem{smith2020standard}
A.~M. Smith, J.~R. Walsh, J. Long, et al.
\newblock Standard machine learning approaches outperform deep representation
learning on phenotype prediction from transcriptomics data.
\newblock \emph{BMC Bioinformatics}, 21:119, 2020.
\newblock doi:10.1186/s12859-020-3427-8.

\bibitem{way2017evaluating}
G.~P. Way and C.~S. Greene.
\newblock Evaluating deep variational autoencoders trained on pan-cancer gene
expression.
\newblock \emph{arXiv preprint arXiv:1711.04828}, 2017.

\bibitem{jiang2025gexbert}
S. Jiang and S. Hassanpour.
\newblock Transformer-based representation learning for robust gene expression
modeling and cancer prognosis.
\newblock \emph{Scientific Reports}, 15:37581, 2025.
\newblock doi:10.1038/s41598-025-14949-2.

\bibitem{xia2022crossstudy}
F. Xia, J.~E. Allen, P. Balaprakash, et al.
\newblock A cross-study analysis of drug response prediction in cancer cell
lines.
\newblock \emph{Briefings in Bioinformatics}, 23(1):bbab356, 2022.
\newblock doi:10.1093/bib/bbab356.

\bibitem{wissel2025survboard}
D. Wissel, N. Janakarajan, A. Grover, E. Toniato, M. Rodr\'iguez Mart\'inez,
and V. Boeva.
\newblock SurvBoard: Standardized benchmarking for multi-omics cancer survival
models.
\newblock \emph{Briefings in Bioinformatics}, 26(5):bbaf521, 2025.
\newblock doi:10.1093/bib/bbaf521.

\bibitem{langfelder2008wgcna}
P. Langfelder and S. Horvath.
\newblock WGCNA: An R package for weighted correlation network analysis.
\newblock \emph{BMC Bioinformatics}, 9:559, 2008.
\newblock doi:10.1186/1471-2105-9-559.

\bibitem{liberzon2011msigdb}
A. Liberzon, A. Subramanian, R. Pinchback, H. Thorvaldsdottir, P. Tamayo, and
J.~P. Mesirov.
\newblock Molecular signatures database (MSigDB) 3.0.
\newblock \emph{Bioinformatics}, 27(12):1739--1740, 2011.
\newblock doi:10.1093/bioinformatics/btr260.

\bibitem{milacic2024reactome}
M. Milacic, D. Beavers, P. Conley, et al.
\newblock The Reactome Pathway Knowledgebase 2024.
\newblock \emph{Nucleic Acids Research}, 52(D1):D672--D678, 2024.
\newblock doi:10.1093/nar/gkad1025.

\bibitem{kanehisa2000kegg}
M. Kanehisa and S. Goto.
\newblock KEGG: Kyoto Encyclopedia of Genes and Genomes.
\newblock \emph{Nucleic Acids Research}, 28(1):27--30, 2000.
\newblock doi:10.1093/nar/28.1.27.

\bibitem{barbie2009systematic}
D.~A. Barbie, P. Tamayo, J.~S. Boehm, et al.
\newblock Systematic RNA interference reveals that oncogenic KRAS-driven
cancers require TBK1.
\newblock \emph{Nature}, 462:108--112, 2009.

\bibitem{jaume2024modeling}
G. Jaume, A. Vaidya, R.~J. Chen, D.~F.~K. Williamson, P.~P. Liang, and
F. Mahmood.
\newblock Modeling dense multimodal interactions between biological pathways
and histology for survival prediction.
\newblock In \emph{CVPR}, pages 11579--11590, 2024.
\newblock doi:10.1109/CVPR52733.2024.01100.

\bibitem{gallaghersyed2026protopathway}
A. Gallagher-Syed, C. Pitzalis, M.~J. Lewis, M.~R. Barnes, and G. Slabaugh.
\newblock ProtoPathway: Biologically structured prototype-pathway fusion for
multimodal cancer survival prediction.
\newblock \emph{arXiv preprint arXiv:2605.21454}, 2026.

\bibitem{brody2022attentive}
S. Brody, U. Alon, and E. Yahav.
\newblock How attentive are graph attention networks?
\newblock In \emph{International Conference on Learning Representations
(ICLR)}, 2022.

\bibitem{hamilton2017inductive}
W.~L. Hamilton, R. Ying, and J. Leskovec.
\newblock Inductive representation learning on large graphs.
\newblock In \emph{Advances in Neural Information Processing Systems},
pages 1024--1034, 2017.

\bibitem{ogris2017binox}
C. Ogris, D. Guala, T. Helleday, and E.~L.~L. Sonnhammer.
\newblock A novel method for crosstalk analysis of biological networks:
improving accuracy of pathway annotation.
\newblock \emph{Nucleic Acids Research}, 45(2):e8, 2017.
\newblock doi:10.1093/nar/gkw849.

\bibitem{schmitt2014funcoup}
T. Schmitt, C. Ogris, and E.~L.~L. Sonnhammer.
\newblock FunCoup 3.0: Database of genome-wide functional coupling networks.
\newblock \emph{Nucleic Acids Research}, 42(D1):D380--D388, 2014.
\newblock doi:10.1093/nar/gkt984.

\bibitem{liu2024pathformer}
X. Liu, Y. Tao, Z. Cai, et al.
\newblock Pathformer: A biological pathway informed transformer for disease
diagnosis and prognosis using multi-omics data.
\newblock \emph{Bioinformatics}, 40(5):btae316, 2024.
\newblock doi:10.1093/bioinformatics/btae316.

\bibitem{lin2023evolutionary}
Z. Lin, H. Akin, R. Rao, et al.
\newblock Evolutionary-scale prediction of atomic-level protein structure with
a language model.
\newblock \emph{Science}, 379(6637):1123--1130, 2023.
\newblock doi:10.1126/science.ade2574.

\bibitem{kang2026bulkformer}
B. Kang, R. Fan, M. Yi, C. Cui, and Q. Cui.
\newblock BulkFormer: A large-scale foundation model for bulk transcriptomes.
\newblock \emph{Cell Systems}, 17(7):101657, 2026.
\newblock doi:10.1016/j.cels.2026.101657.

\bibitem{choromanski2021rethinking}
K. Choromanski, V. Likhosherstov, D. Dohan, et al.
\newblock Rethinking attention with performers.
\newblock In \emph{International Conference on Learning Representations
(ICLR)}, 2021.

\end{thebibliography}

\end{document}
